\section{红黑树后端：\texttt{vma\_rbtree.c}}
\label{sec:vma-rbtree}

排序数组已经满足了早期内核的需求，但当进程管理引入后，VMA 数量会急剧增长。
红黑树在所有单次操作上都保持 $O(\log n)$，是 Linux 内核历史上长期使用的方案
（Linux 2.4 $\sim$ 6.0，整整 20 年）。

\begin{table}[H]
\centering
\begin{tabular}{lcccc}
\toprule
\textbf{操作} & \textbf{链表} & \textbf{排序数组} & \textbf{红黑树} & \textbf{Maple Tree} \\
\midrule
\texttt{find}(addr)   & $O(n)$      & $O(\log n)$  & $O(\log n)$  & $O(\log_B n)$ \\
\texttt{add}          & $O(n)$      & $O(n)$       & $O(\log n)$  & $O(\log_B n)$ \\
\texttt{remove}       & $O(n)$      & $O(n)$       & $O(\log n)$  & $O(\log_B n)$ \\
\texttt{dump}（遍历） & $O(n)$      & $O(n)$       & $O(n)$       & $O(n)$        \\
\bottomrule
\end{tabular}
\caption{不同数据结构的 VMA 操作复杂度对比（$B$ = B-tree 扇出度）}
\label{tab:vma-complexity}
\end{table}

\subsection{红黑树的 5 条性质}

红黑树是一种自平衡二叉搜索树，每个节点带一个颜色位（红/黑），
且满足以下 5 条性质：

\begin{enumerate}
  \item 每个节点要么是红色，要么是黑色。
  \item 根节点是黑色。
  \item 所有叶子节点（NIL 哨兵）是黑色。
  \item 红色节点的两个子节点必须是黑色（\emph{不能有连续的红节点}）。
  \item 从任意节点到其所有后代叶子节点的路径上，黑色节点数相同（\emph{黑高相等}）。
\end{enumerate}

性质 4 + 5 共同保证了树的高度不超过 $2\log_2(n+1)$，
从而使查找、插入、删除都是 $O(\log n)$。

\subsection{内部节点结构}

每个红黑树节点嵌入一个 \texttt{vm\_area\_t} 公开视图，外加颜色位和三个指针：

\codefile{src/kernel/mm/vma\_rbtree.c}
\begin{lstlisting}[style=cstyle]
#define RB_MAX_NODES 256
#define RB_BLACK 0
#define RB_RED   1

typedef struct rb_node {
    vm_area_t       area;   /* 公开视图（第一个字段） */
    int             color;  /* RB_BLACK or RB_RED */
    struct rb_node* parent;
    struct rb_node* left;
    struct rb_node* right;
} rb_node_t;
\end{lstlisting}

树使用一个共享的 NIL 哨兵节点（颜色始终为黑）替代 NULL，
消除了每次访问子指针前的判空检查，极大简化旋转和修复代码。

\subsection{旋转操作}

当插入或删除破坏了红黑性质时，需要通过\term{旋转}（rotation）恢复平衡。
旋转分为左旋和右旋，它们互为镜像操作：

\begin{figure}[H]
  \centering
  % figures/ch11/fig07-rbtree-rotate.tex — auto-extracted from chapter source
\begin{tikzpicture}[
  node distance=1cm and 1.2cm,
  every node/.style={draw, circle, minimum size=0.7cm, font=\small},
  tri/.style={draw, isosceles triangle, shape border rotate=90,
              minimum height=0.5cm, minimum width=0.6cm, inner sep=0pt,
              font=\scriptsize},
  arr/.style={-{Stealth[length=5pt]}, very thick, blue!60!black},
]
  \node[fill=red!20] (x) at (0, 0) {$x$};
  \node[tri] (a) at (-1, -1.2) {$\alpha$};
  \node[fill=red!20] (y) at (1, -1.2) {$y$};
  \node[tri] (b) at (0.3, -2.4) {$\beta$};
  \node[tri] (c) at (1.7, -2.4) {$\gamma$};
  \draw (x) -- (a);
  \draw (x) -- (y);
  \draw (y) -- (b);
  \draw (y) -- (c);

  \node[draw=none, font=\Large\bfseries] at (3.5, -1) {$\Rightarrow$};
  \node[draw=none, font=\small] at (3.5, -1.8) {左旋 $x$};

  \node[fill=red!20] (y2) at (6.5, 0) {$y$};
  \node[fill=red!20] (x2) at (5.5, -1.2) {$x$};
  \node[tri] (c2) at (7.5, -1.2) {$\gamma$};
  \node[tri] (a2) at (4.8, -2.4) {$\alpha$};
  \node[tri] (b2) at (6.2, -2.4) {$\beta$};
  \draw (y2) -- (x2);
  \draw (y2) -- (c2);
  \draw (x2) -- (a2);
  \draw (x2) -- (b2);
\end{tikzpicture}

  \caption{红黑树左旋示意图（右旋为其镜像）}
  \label{fig:rbtree-rotate}
\end{figure}

旋转只改变指针指向，不改变 BST 的中序遍历顺序——
因此旋转前后所有 VMA 的地址排序仍然正确。

\codefile{src/kernel/mm/vma\_rbtree.c}
\begin{lstlisting}[style=cstyle]
static void rotate_left(rb_node_t* x)
{
    rb_node_t* y = x->right;
    x->right = y->left;
    if (y->left != NIL)
        y->left->parent = x;
    y->parent = x->parent;
    if (x->parent == NIL)
        g_root = y;
    else if (x == x->parent->left)
        x->parent->left = y;
    else
        x->parent->right = y;
    y->left = x;
    x->parent = y;
}
\end{lstlisting}

右旋 (\texttt{rotate\_right}) 是完全对称的镜像操作，交换 left/right 即可。

\subsection{插入与修复}

新节点默认为红色，BST 插入后可能违反性质 4（连续红节点）。
修复分 3 种情况（加镜像共 6 种）：

\begin{enumerate}
  \item \textbf{叔节点为红}——父和叔变黑，祖父变红，向上递归。最坏 $O(\log n)$ 次重着色。
  \item \textbf{叔节点为黑 + z 是内侧孩子}——旋转父节点，转为情况 3。
  \item \textbf{叔节点为黑 + z 是外侧孩子}——旋转祖父 + 重着色，修复完成。
\end{enumerate}

最坏情况：$O(\log n)$ 次重着色 + 最多 2 次旋转。

\subsection{删除与修复}

删除一个黑色节点会破坏性质 5（黑高不等），修复分 4 种情况（加镜像共 8 种）。
这是红黑树最复杂的部分，但最坏也只需 $O(\log n)$ 次重着色 + 最多 3 次旋转。

\subsection{区间查找}

BST 按 \texttt{start} 排序，但 VMA 是区间 \texttt{[start, end)}。
查找算法从根下降，每步做三路判断：

\codefile{src/kernel/mm/vma\_rbtree.c}
\begin{lstlisting}[style=cstyle]
static const vm_area_t* rb_find(uint32_t addr)
{
    rb_node_t* cur = g_root;
    while (cur != NIL) {
        if (addr < cur->area.start)
            cur = cur->left;
        else if (addr >= cur->area.end)
            cur = cur->right;
        else
            return &cur->area;   /* 命中 */
    }
    return NULL;
}
\end{lstlisting}

\subsection{中序遍历 dump}

利用 parent 指针做迭代中序遍历（不递归，避免栈溢出风险），
天然按地址升序输出。

\begin{thinkbox}
如果 VMA 后端改用简单的有序链表（插入时遍历找位置），
在 VMA 数量为 $n=200$ 时，\texttt{find} 的最坏情况要遍历多少个节点？
红黑树呢？（提示：$\lceil \log_2 201 \rceil = 8$）

在什么场景下链表反而比红黑树更合适？
（提示：VMA 数量极少且变化频繁——旋转的常数开销）
\end{thinkbox}

% ============================================================================

